% !TeX root = ../thesisv2.tex

\chapter{结论与展望}
\label{cha:v2-conclusion}

\section{结论}
本文围绕高分辨率遥感影像道路提取任务，在当前本地全部可复核材料基础上，完成了实验工程重建、统一 baseline 建立、主线模型设计、baseline 改进分析以及多模型整图对照实验。结合前文理论分析、模型设计和实验结果，可以形成以下结论。

第一，本文建立了完整、统一、可复用的遥感道路提取实验工程。该工程覆盖数据整理、离线 patch 生成、模型训练、整图推理、整图评估、结果可视化和论文支撑等主要环节，解决了原始材料中代码缺失、训练链路不闭合和结果难以统一比较的问题。通过离线 patch 数据集构建、多进程数据加载和日志同步优化，基线模型单 epoch 训练耗时由 409.3 s 降至 113.5 s，训练效率提升约 3.61 倍，为后续多模型完整实验提供了必要条件。

第二，统一 baseline U-Net 构成了全文模型比较的公共参照。基于 Batch Normalization 和双线性上采样的 baseline U-Net 在 Massachusetts Roads Dataset 测试集上达到 IoU 0.645964、F1 0.782378、Precision 0.787734 和 Recall 0.780151。该模型参数量仅为 7.85M，在较低参数规模下表现出稳定、均衡且可复用的综合性能。因此，baseline 不仅是后续模型比较的起点，也是本文判断结构改动是否有效的核心标准。

第三，DLGU-Net 构成了本文的主线模型。该模型在四级跳跃连接中引入轻量级双线性注意力模块 DLAM，通过通道重标定和空间方向增强对浅层特征进行筛选与增强。实验结果表明，DLGU-Net 在约 8.30M 参数规模下达到 Test IoU 0.641934 和 Test F1 0.779244，并在 Precision 指标上表现出较为明显的背景抑制能力。结合与 Attention U-Net、UNet++ 的结构对照可以确认，跳跃连接增强是一个具有明确研究意义的方向，而 DLGU-Net 为这一方向提供了结构清晰、参数开销可控、训练过程稳定的实现方案。

第四，围绕 baseline 的直接增强路线中，空洞卷积增强 baseline 取得了最具代表性的正向结果。该模型仅在 bottleneck 位置引入并行空洞卷积分支，其余结构保持 baseline 风格不变，在测试集上达到 IoU 0.646839 和 F1 0.782933，相比 baseline 分别提升 0.000875 和 0.000555，并在 Precision 上提升 0.023358。由此可见，在当前任务条件下，于高层语义特征位置增强多尺度上下文表达，是围绕 baseline 获得稳定增益的有效方式。

第五，扩展探索分支显著丰富了全文的结构比较维度。Vanilla U-Net 在全体模型中取得最高测试集 IoU 0.647187 和 F1 0.783246，说明在当前数据规模、统一训练设置和整图评估协议下，原始风格卷积块与可学习上采样具有较好的任务适配性。Residual Vanilla U-Net、ResNet34 U-Net 和 DDU-Net 则分别从残差连接、预训练编码器和重型上下文结构角度补充了结果矩阵。通过这些对照，本文进一步证明：模型表现并不简单取决于参数量大小，而更取决于结构改动是否与道路提取任务特征相匹配。

第六，本文形成了层次清晰的模型角色划分。统一 baseline 用于提供稳定参照；DLGU-Net 作为主线模型，用于验证轻量级跳跃连接增强方向；Dilated Baseline U-Net 作为 baseline 直接改进路线中的代表结果，用于支撑“围绕 baseline 取得小幅提升”的论文主线；Vanilla U-Net 则作为扩展探索分支中的最佳结果，体现本文实验广度和结构搜索深度。通过这种层次化组织，全文既保持了明确主线，也完整呈现了全部实际工作量。

\subsection*{工程层面结论}
从工程层面看，本文最重要的结果在于完成了原始实验工程向统一可复核平台的重建。数据目录标准化、配置驱动训练、模型工厂、整图评估模块和自动可视化导出等功能使不同模型能够在同一工作流下完成训练和测试。对于毕业论文而言，这种工程完整性保证了结论来源清晰，也为后续继续扩展模型和补充实验保留了稳定基础。

\subsection*{方法层面结论}
从方法层面看，本文验证了两类结构增强路线在当前任务中的作用特征。第一类是跳跃连接增强路线，以 DLGU-Net、Attention U-Net 和 UNet++ 为代表，重点解决浅层特征筛选与融合问题。第二类是高层上下文增强路线，以空洞卷积增强 baseline 为代表，重点解决瓶颈层多尺度上下文表达问题。实验结果说明，二者都具有研究价值，但在当前数据条件下，高层上下文增强更容易转化为直接的测试集精度收益。

\subsection*{实验层面结论}
从实验层面看，统一整图评估协议是本文结论可信的重要前提。训练阶段使用 patch 输入，验证和测试阶段使用整图切块恢复与整图统计，使全部模型能够在相同口径下进行比较。最终结果表明，Vanilla U-Net、Dilated Baseline U-Net、Baseline U-Net、UNet++ 和 DLGU-Net 形成了当前材料下的主要结果梯队，而更复杂的重型结构未在当前任务中体现出明显优势。这一结论说明，针对道路提取任务的结构适配性比单纯增加模型复杂度更重要。

\section{展望}
在本文工作基础上，后续研究仍有多个方向值得继续深入。

第一，围绕 DLGU-Net 可继续开展更细粒度的组件级研究。当前主线模型已经完成完整实现和统一评估，后续可进一步分别考察通道分支与空间分支的独立贡献、不同条形卷积核尺寸的影响、不同 reduction 比例对参数开销与性能的平衡效果，以及 DLAM 模块插入层级变化对最终结果的影响。这类研究有助于将当前主线模型从“完整模块验证”推进到“组件机制分析”。

第二，围绕空洞卷积增强 baseline 可继续深化多尺度上下文建模策略。现有模型使用膨胀率 1、2、4、8 的并行分支，后续可以进一步分析不同膨胀率组合、分支融合方式、bottleneck 前后插入位置和上下文模块深度变化对精度与参数量的影响。由于该路线已经在统一测试集上形成正向收益，因此具有较高继续优化价值。

第三，后续可在其他公开遥感道路数据集上补充跨数据集泛化实验。当前完整可复核实验主要集中在 Massachusetts Roads Dataset 上，统一工程已经具备迁移到其他道路提取数据集的能力。若进一步补充 DeepGlobe 等数据集上的训练和测试，可更全面分析不同模型在复杂场景、不同标注风格和不同地理分布下的适配性。

第四，可将道路拓扑约束与后处理机制引入当前实验平台。道路提取任务的最终目标不仅是像素分类准确，还包括道路网络的结构连通与几何完整。后续可在当前分割模型基础上结合图结构建模、中心线恢复、连通性约束或拓扑一致性损失，使模型输出更加贴近真实道路网络的组织方式。

第五，可围绕轻量化部署继续开展模型压缩与推理优化研究。本文已经通过工程优化显著提高训练效率，但在模型部署层面仍可结合剪枝、蒸馏、轻量 backbone 替换和推理图优化等技术，构建更适合边缘设备、移动终端或星载平台的道路提取方案。对于需要实时响应或大规模批处理的遥感应用，这一方向具有明确工程价值。

第六，可进一步加强论文工程与研究工程的一体化组织。本文已经将训练结果、整图评估、可视化图像和论文图表组织在同一工程中。后续可以继续完善自动化图表生成、自动摘要更新和实验记录归档机制，使论文写作与实验执行之间形成更紧密的闭环，进一步提升结果复核效率和研究积累效率。

总体来看，本文已经形成从数据整理、工程重建、模型实现、统一评估到论文组织的完整工作链条。该链条既支撑了当前毕业论文的完成，也为后续继续开展遥感道路提取结构设计与工程研究提供了稳定起点。
